%%%%%%%% MLSys 2026 submission (anonymous draft) %%%%%%%%%%%%%%%%%%%%%%%%%%%%
\documentclass{article}

\usepackage{microtype}
\usepackage{graphicx}
\usepackage{booktabs}
\usepackage{amsmath}
\usepackage{multirow}
\usepackage{array}
\usepackage{hyperref}
\newcommand{\theHalgorithm}{\arabic{algorithm}}

% Blind submission: the style file removes author information automatically.
\usepackage{mlsys2025}
\mlsystitlerunning{Joint Precision Budgeting for Hybrid Linear-Attention Serving}

\begin{document}

\twocolumn[
\mlsystitle{Joint Precision Budgeting Across Attention KV and Recurrent State in Hybrid Linear-Attention Serving}

\begin{mlsysauthorlist}
\mlsysauthor{Anonymous}{}
\end{mlsysauthorlist}
\mlsysaffiliation{}{Anonymous Institution}
\mlsyscorrespondingauthor{Anonymous Author}{anon.email@domain.com}

\vskip 0.3in

\begin{abstract}
Hybrid linear-attention models pair a small set of GQA attention layers with
interleaved Gated Delta Network (GDN) layers, so serving memory is split
between a growing, quantizable attention KV cache and a fixed per-sequence
recurrent state that no KV quantizer touches. Existing systems optimize the
KV bit-width and leave state precision fixed. We treat state precision as a
second budget dimension and implement a selector that maps quality, memory,
concurrency, and latency constraints to executable precision configurations.
On Qwen3.5-2B/9B served with vLLM on one RTX 5090, a reproducible 112-cell
allocator matrix shows higher capacity in all 52 paired fp32/bf16 cells, with
a median gain of 15.44\%. The idealized capacity model has 1.81\% median
absolute residual (range $-3.66\%$ to $+13.21\%$); it is a prediction, not a
bound. Quality remains task-dependent: GSM8K loses 1.0 percentage point at
2B, whereas five preregistered no-think RULER cells show exact observed
agreement across three dataset seeds, with all 30 accuracy values and sample
outputs reproduced exactly. This equality does not establish equivalence.
The selector reproducibly maps three preregistered budgets to full,
state-only, and joint precision, but only for one 2B/4K Random-workload slice.
Controlled mechanism contrasts reproduce throughput, not tail latency. A
broader four-configuration serving matrix completes 144/144 cells per attempt
but fails its continuous-goodput stability gate (183/720 comparisons outside
10\% tolerance). We therefore claim a verified capacity effect and scoped
policy execution, while reporting serving instability rather than a universal
performance benefit.
\end{abstract}
]

\printAffiliationsAndNotice{}

\section{Introduction}
\label{sec:intro}

KV caches are the dominant memory consumer in LLM serving, and KV
quantization is the standard lever for fitting more context and concurrency
into a fixed GPU budget~\cite{kivi,kvquant,minikv}. The field has largely
been built on attention-only transformers, in which every layer contributes a
growing, homogeneous, quantizable cache.

Hybrid linear-attention models break that assumption. Qwen3.5, the subject of
this paper, interleaves 18 (2B) or 24 (9B) Gated Delta Network (GDN)
layers~\cite{gateddeltanet} with 6 (2B) or 8 (9B) GQA attention layers. The
GDN layers keep no per-token KV; instead each sequence carries a fixed
recurrent state that is allocated from the same GPU memory pool as the
attention KV. Consequently the serving memory footprint of a hybrid model is
two-dimensional: a per-token component that KV quantization can shrink, and a
per-sequence state component whose size depends on its own data type. State
precision is a deployment switch in vLLM (\texttt{mamba\_ssm\_cache\_dtype}),
but existing budget analyses, including prior work on state compression
itself~\cite{replayssm,vllmpr43518}, do not treat it as a dimension that
should be jointly allocated with the KV bit-width.

This paper is a system study with four contributions. (a) A capacity model
that separates concurrent sequences $N(L)=M/(A L+G)$ from total token
capacity $T(L)=L N(L)$, plus a deterministic probe validated by two complete
112-cell attempts. (b) An executable joint-precision selector that filters
full, KV-only, state-only, and joint candidates by quality, memory,
concurrency, and tail-latency constraints, then maximizes a robust SLO-goodput
lower bound and fails closed on unmeasured strata. (c) A paired quality map
across perplexity, GSM8K, and no-think RULER, with statistical resolution and
reproduction status stated per endpoint. (d) Controlled mechanism contrasts
and a four-configuration serving matrix whose partial and failed
reproducibility results are retained rather than converted into a performance
headline.
The target regime is short-context, high-concurrency, memory-bound serving,
where the fixed state is a large fraction of the budget.

The headline result is a verified capacity effect. Across the frozen matrix
of model size, KV dtype, context length, and memory utilization, all 52
fp32/bf16 pairs favor bf16 state and the median gain is 15.44\%. The clean
R2-to-R3 comparison reproduces all 112 cells within the preregistered
tolerances. The architecture-derived model has a median absolute residual of
1.81\%, with signed residuals from $-3.66\%$ to $+13.21\%$; page and block
rounding therefore make it an idealized predictor rather than a lower or upper
bound. In a representative 2B/4K cell, the int4-over-fp16 KV capacity ratio
rises from 2.245 to 2.675 when state is also bf16.

The quality evidence is bounded. Perplexity changes are within
$[-0.013, +0.007]$ (C4) and $[-0.045, +0.058]$ (PG19) with all confidence
intervals containing zero. GSM8K (200 seeded items, 9 seeds) shows a
significant 1.0-point regression at 2B state-bf16 ($p=0.025$), which we
report together with its pre-registered minimum detectable effect of 1.16
points and 67.5\% observed power; the 9B model shows no regression. In the
no-think RULER repair, every tested fp32/bf16 pair has a delta of 0.00 across
three dataset seeds, and an independent temporal attempt reproduces all 30
accuracy values and sample outputs exactly. These are observed agreements,
not equivalence or non-inferiority results. The controller's three mappings
also reproduce within their frozen tolerances. In contrast, tail-latency
mechanism metrics and the broader four-configuration goodput matrix do not
pass their respective stability gates, so we make no general serving-speed
claim.

\section{Background and Related Work}
\label{sec:related}

\textbf{KV cache quantization and eviction.} KIVI~\cite{kivi} and
KVQuant~\cite{kvquant} established 2--4-bit per-token-head KV quantization;
MiniKV~\cite{minikv} co-designs 2-bit layer-discriminative KV quantization
with the vLLM system, while TurboQuant~\cite{turboquant} applies online
vector quantization to the KV cache at 2--4 bits per entry in vLLM. Neither
work treats the recurrent-state dtype as an allocatable dimension. Eviction
and hybrid policies
(H2O~\cite{h2o}, StreamingLLM~\cite{streamingllm}, SnapKV~\cite{snapkv})
release memory by dropping tokens rather than reducing precision. All of
these target attention-only models. On a hybrid model they compress only the
attention layers; the GDN state remains untouched. Eviction and
offloading therefore operate on a different axis from the state-dtype switch:
they change which tokens or layers stay in memory, whereas state-bf16 changes
the bytes per recurrent-state element. Our formal matrix holds eviction and
offloading fixed, so the two levers are not mixed in the measured deltas.
Recent system-level work
also co-optimizes eviction with quantization within the attention-KV budget
(HqeKV~\cite{hqekv}, RDKV~\cite{rdkv}, ARKV~\cite{arkv}).

\textbf{Hybrid linear-attention architectures.} Mamba~\cite{mamba} and
Mamba2~\cite{mamba2} established selective state-space layers; Gated Delta
Networks~\cite{gateddeltanet} make the recurrence a delta rule. Production
hybrids (Jamba~\cite{jamba}, RecurrentGemma~\cite{recurrentgemma},
Zamba~\cite{zamba}, Qwen3.5~\cite{qwen35}) interleave such layers with a small
number of full-attention layers. RecurrentGemma builds on
Griffin~\cite{griffin}; the GDN recurrence continues the delta-rule lineage
of DeltaNet~\cite{deltanet}. Serving them requires a per-sequence state
in addition to the attention KV; vLLM's PagedAttention memory
manager~\cite{vllm} now allocates both from the same GPU pool.

\textbf{State precision and compression.} Quantizing SSM states has been
explored in the low-precision inference literature: ReplaySSM proposes
caching inputs instead of state, and vLLM exposes a configurable
state-dtype path~\cite{replayssm,vllmpr22196}; FP8 state checkpointing is
under development in the FlashInfer SSU path~\cite{vllmpr43518}. Recurrent
state caching also appears in inference frameworks such as
FlashInfer~\cite{flashinfer}.
Weight-and-activation quantization of SSM layers
(Quamba~\cite{quamba}, MambaQuant~\cite{mambaquant}) targets the compute
path with W8A8 or W4A8 schemes and does not vary the per-sequence state
dtype that serving allocators budget, which is the dimension we vary.
QPruningKV~\cite{qpruningkv} characterized the token-precision trade-off
under a fixed byte budget for attention models. State compression itself is
prior art. The gap is joint budget accounting: serving systems expose state
dtype without analyzing it as a precision dimension that multiplies with
the KV bit-width decision, and without validating the resulting
capacity-quality-serving trade-off with disclosed statistical resolution and
a second formal run.

\section{Precision Budget Model}
\label{sec:model}

Let $L$ be the context length, $A_f$ and $A_q$ the per-token attention-KV
bytes at fp16 and int4 KV respectively, and $G$ the per-sequence state bytes
at the chosen state dtype. With $M$ bytes available for KV plus state, the
maximum number of concurrent length-$L$ sequences and the corresponding
total token capacity are
\begin{equation}
  N(L) = \frac{M}{A\cdot L + G}, \qquad
  T(L) = L N(L) = \frac{L\cdot M}{A\cdot L + G},
  \label{eq:capacity}
\end{equation}
where $A\in\{A_f,A_q\}$ and $G$ depends on the state dtype; the denominator
is the per-sequence footprint of a length-$L$ request. The capacity ratio from
switching the state fp32 to bf16 is
\begin{equation}
  r_{state}(L)=\frac{A L + G_{fp32}}{A L + G_{bf16}},
  \label{eq:rstate}
\end{equation}
which decreases with $L$: state dominates at short context, KV dominates at
long context. The corresponding KV ratio at a fixed state dtype is
$r_{kv}(L)=(A_f L+G)/(A_q L+G)$. The parameters $A_q$, $A_f$, and $G$ are
derived from the model configuration and the vLLM page layout (per-layer
state shapes and padded block sizes), independent of the probed token
counts, so the comparison of Eq.~\eqref{eq:rstate} against measurements is
not circular. Details of the derivation and the padded block accounting are
in the appendix.

For deployment stratum $x$ and candidate precision configuration $q$, the
selector implements
\begin{equation}
q^*(x)=\arg\max_{q\in\mathcal{F}(x)}
\underline{\mathrm{Goodput}}(q,x),
\label{eq:selector}
\end{equation}
where $\mathcal{F}(x)$ contains only measured candidates satisfying
$\underline{\Delta Q_t}\ge-\epsilon_t$ for every requested task,
$\overline{\mathrm{TTFT}}_{95}\le\tau_f$,
$\overline{\mathrm{TPOT}}_{95}\le\tau_p$, cache bytes $\le M$, and
$N_q(L)\ge N_{min}$. Underlines and overlines denote the stored lower and
upper confidence bounds. Missing or ambiguous profile rows yield no feasible
candidate rather than interpolation or extrapolation.

\section{Evaluation Setup}
\label{sec:setup}

\textbf{Platform and engine.} All experiments use Qwen3.5-2B or Qwen3.5-9B
on a single RTX 5090 (32 GB). Earlier serving and quality matrices use a
frozen vLLM worktree at commit \texttt{55f4768}; the preregistered capacity,
repaired RULER, controller, mechanism, and four-configuration packages use
the later frozen worktree at commit \texttt{e2fa285}. Memory utilization is a
protocol-level parameter rather than a global constant: 0.85 for the earlier
matrices and RULER, 0.70--0.90 in the capacity grid (plus 0.85 controls), and
0.80 for the controller and follow-up serving packages. Each package archives
its resolved cache dtypes and protocol fields, and repetitions are compared
only within the corresponding frozen package.

\textbf{Capacity probes.} We use the engine's maximum-token probe with
1K--32K contexts, uniform int4 KV (\texttt{int4\_per\_token\_head}) or fp16
KV, state fp32/bf16, and the preregistered utilization grid. Each cell is a
single deterministic allocator run; each attempt records the resolved dtype,
block counts, and serving configuration. Two clean attempts each complete
112/112 cells (52 fp32/bf16 core pairs plus 8 fp16-state controls), and Gate~4
compares per-cell tokens, directions, and model-by-KV median gains.
Table~\ref{tab:capacity} shows an independently verified fixed-utilization
4K/16K slice used for deployment interpretation.

\textbf{Quality.} Perplexity is measured with a chunk-level harness
(write-back rounding at chunk boundaries, chunk 128) on C4 and PG19, 5
sequences, 3 dataset seeds; we report in Section~\ref{sec:limits} why this
harness is an approximation. Reasoning quality uses the vLLM kernel path:
GSM8K~\cite{gsm8k} with 200 seeded items per seed, greedy, no
chain-of-thought, 9 seeds; RULER~\cite{ruler} on five preregistered
model-task-length cells, 20 samples per cell, 3 dataset seeds, fixed engine
seed, thinking disabled, and a 256-token generation cap. All comparisons are
paired across dataset seeds. For RULER, all paired deltas are identical zeros,
so the zero-width descriptive intervals denote observed equality only.

\textbf{Joint-precision controller.} Each deployment request specifies the
model and context, workload and offered load, memory and concurrency budgets,
TTFT/TPOT SLOs, and a task-level quality floor. The controller requires an
exact measured profile row for every candidate and does not interpolate across
strata. It rejects candidates whose quality lower bound, cache bytes,
concurrency, or latency upper bounds violate the request, maximizes the
remaining SLO-goodput lower confidence bound, and maps the selected KV/state
dtypes to a cold-restart vLLM command. We test three frozen budgets on
Qwen3.5-2B at 4K, Random workload, TP=1, with seeds 11, 23, and 47. The parent
and temporal reproduction each contain 9 samples and 18,000 measurement
requests, with failures retained in the SLO denominator.

\textbf{Serving.} Protocol v3: piecewise constant Poisson arrivals,
120 s warmup, 60 s (300 s for boundary confirmation) measurement windows,
TTFT thresholds 250--3000 ms, TPOT 200 ms, goodput/offered $\ge 0.95$
required for a sustainable boundary, failures counted in the denominator,
3 seeds. Workloads: Random60 (60-token requests) and ShareGPT300 (300-token
traces~\cite{sharegpt}). The formal matrix was followed by a second formal run with the same
contracts and seeds (run-stability evidence, not an independent sample); both are committed as atomic
artifacts~\cite{mlsys2026ae}.

\textbf{Mechanism and four-configuration contracts.} Mechanism isolation
crosses full and joint precision with fixed block count, fixed cache bytes,
and fixed concurrency (3 seeds; 18 cells and 5,400 requests per attempt).
The broader serving matrix crosses full, KV-only, state-only, and joint
precision with Random and ShareGPT workloads (3 seeds; 144 cells and 320,400
requests per attempt). Its frozen Gate~4 compares all 720
cell-threshold goodput values at a strict 10\% run-stability tolerance.

\begin{table}[t]
\centering
\footnotesize
\caption{\textbf{Capacity ratio $r_{state}$ from bf16 state.} Measured token
capacity ratio and architecture-derived prediction; signed gap is
(measured$-$predicted)/predicted. This 4K/16K slice illustrates allocator
rounding; the full 52-pair matrix has residuals on both sides of zero, so the
model is not a bound.}
\label{tab:capacity}
\setlength{\tabcolsep}{4pt}
\begin{tabular}{lccccl}
\toprule
KV & Model & $L$ & fp32 tok & bf16 tok & gap\\
\midrule
int4 & 2B & 4096 & 2,692,710 & 3,703,954 & $-2.37\%$\\
int4 & 2B & 16384 & 4,895,837 & 5,458,458 & $-3.24\%$\\
int4 & 9B & 4096 & 315,392 & 443,538 & $-0.18\%$\\
int4 & 9B & 16384 & 573,440 & 653,635 & $-1.07\%$\\
fp16 & 2B & 4096 & 1,199,383 & 1,384,448 & $+2.86\%$\\
fp16 & 2B & 16384 & 1,552,143 & 1,661,337 & $+3.53\%$\\
fp16 & 9B & 4096 & 144,104 & 161,899 & $-2.83\%$\\
\bottomrule
\end{tabular}
\end{table}

\section{Results}
\label{sec:results}

\subsection{Capacity}
\label{sec:capacity}

Table~\ref{tab:capacity} and Figure~\ref{fig:capacity} report the capacity
measurements. Under int4 KV, bf16 state raises capacity by 37.6\% (2B, 4K),
40.6\% (9B, 4K), 11.5\% (2B, 16K), and 14.0\% (9B, 16K). The fp16-KV column
shows smaller gains (7.0--15.4\% at 4K) with mixed prediction signs. The
discrete page-block rounding is the source of the residual: fp32/bf16 block
sizes change with layout (e.g., int4 2B: 2064/1072 tokens per block; fp16:
544/288). Across the complete 112-cell attempt, all 52 core pairs favor bf16
state; the median capacity gain is 15.44\%. The model's median absolute
residual is 1.81\%, with signed residuals from $-3.66\%$ to $+13.21\%$.
Because residuals cross zero, we use the model as an idealized predictor and
make no bound claim. The clean second attempt reproduces all 112 cells and 52
directions within the frozen Gate~4 tolerances (maximum per-cell token
difference 1.42\%).
Figure~\ref{fig:blocks} shows the mechanism: bf16 state halves the per-block
state footprint, which doubles the number of allocatable GPU blocks, and the
measured ratio is the integer-arithmetic outcome of that re-partitioning.

Figure~\ref{fig:capacity}b also shows the compound KV ratio: with fp32 state,
int4 KV yields 2.2451x over fp16 KV at 2B/4K; with bf16 state the same
comparison yields 2.6754x. Both components draw from the same GPU pool but
enter the per-sequence footprint as the additive terms $A L$ and $G$; changing
either term changes the ratio of feasible capacities.
In deployment terms, the fp32 GDN state at 2B is 18.63 MiB per sequence
(9.32 MiB at bf16). At 4K under int4 KV, the measured total-token ceilings
of 2,692,710 and 3,703,954 tokens correspond to 657 and 904 concurrent
length-4096 sequences on the same GPU budget, so the state switch adds
roughly 247 sequence slots without changing the KV bytes themselves.

\begin{figure*}[t]
\centering
\includegraphics[width=\textwidth]{figures/vector_redesign/fig5_block_granularity.pdf}
\caption{\textbf{Block granularity under fp32 vs.\ bf16 state.}
(a) Tokens per GPU block and (b) number of allocated GPU blocks for the seven
capacity cells. bf16 halves the per-block state footprint; the discrete
re-partitioning of blocks explains the residual between measured and modeled
$r_{state}$ in Table~\ref{tab:capacity}.}
\label{fig:blocks}
\end{figure*}

\begin{figure*}[t]
\centering
\includegraphics[width=\textwidth]{figures/vector_redesign/fig1_capacity.pdf}
\caption{\textbf{Precision budget for hybrid serving.} (a) GDN recurrent
state (fixed per sequence) and attention KV (growing with context) both draw
from the same GPU pool; both bit-width (KV) and state precision are
allocatable. (b) Measured vs.\ predicted $r_{state}$ for the int4/fp16 KV
columns (7 cells; gap labels are signed percent error). int4 column is a
representative slice; the full matrix contains residuals of both signs.}
\label{fig:capacity}
\end{figure*}

\subsection{Quality: perplexity and reasoning}
\label{sec:quality}

\textbf{Perplexity.} Under int4 KV, the bf16-state marginal perplexity change
is $-0.0029$ [$-0.0129$, $+0.0072$] on C4 and $+0.0065$ [$-0.0447$,
$+0.0578$] on PG19 (3 seeds, paired 95\% CI). Both intervals contain zero,
and the magnitude is the same as the state change under fp16 KV
($-0.0003$/$+0.0004$), so stacking state-bf16 onto int4 KV introduces no
measurable additional perplexity cost. The harness is a chunk-level
approximation (Section~\ref{sec:limits}); these numbers are supporting
evidence for the kernel-level claim, not a substitute for it.
Figure~\ref{fig:harness} quantifies both statements: chunk-boundary rounding
moves PPL by roughly 87\% when the chunk size is reduced to one, while the
marginal state cost under int4 KV stays within 0.007 of the fp16-KV cost on
both corpora.

\textbf{GSM8K.} Figure~\ref{fig:gsm8k} summarizes the 9-seed paired deltas.
At 2B, bf16 state loses 1.00 point [$-$1.71, $-$0.29], $p=0.025$, with a
pre-registered minimum detectable effect of 1.16 points and 67.5\% observed
power; this is a significant regression whose effect size is below the
pre-registered MDE, and we report both. Across the nine paired dataset seeds
the regression has paired Cohen's $d=-0.92$; at 9B the paired $d=+0.54$, but
the confidence interval contains zero, so we do not interpret it as a gain.
Uniform int4 KV loses 2.72 points
[$-$4.20, $-$1.24], $p=0.007$, power 88.3\%. Stacking state-bf16 onto int4 KV
shows a marginal $+1.17$ points [$-0.33$, $+2.66$] relative to int4 alone:
no additional regression. At 9B, bf16 state gains 0.33 points [$-0.07$,
$+0.73$], $p=0.141$: no detectable regression. We do not aggregate the 2B and
9B effects. Figure~\ref{fig:gsm8k-seeds} shows the same data at seed
resolution: the nine thin lines pair each dataset seed across allocations,
and the between-seed spread is what the confidence intervals in
Figure~\ref{fig:gsm8k} summarize.

\begin{figure}[t]
\centering
\includegraphics[width=\linewidth]{figures/vector_redesign/fig2_gsm8k.pdf}
\caption{\textbf{GSM8K paired deltas.} 200 seeded items, greedy, no
chain-of-thought, 9 dataset seeds; points and 95\% paired CIs. Red intervals
exclude zero. The 2B state and int4 regressions are significant; the 2B state
effect is below the pre-registered MDE (1.16 pp) with 67.5\% power.}
\label{fig:gsm8k}
\end{figure}

\textbf{RULER.} Figure~\ref{fig:ruler} reports the five cells (FWE and
NIAH-multiquery, 4K/8K) that the earlier KV-quantization screen had singled
out; the selection predates the state-dtype grid and is not a post hoc
choice from these outcomes. Thinking is disabled to remove truncation of the
model's hidden reasoning from the 256-token answer budget. Across the three
paired dataset seeds, fp32 and bf16 state produce identical accuracy in every
cell: 91.11 and 98.33 for 2B FWE at 4K and 8K, and 100.00 for the three 9B
cells. Thus all five bf16-minus-fp32 deltas are 0.00. An independent temporal
attempt reproduces all 30 allocation-by-seed accuracy values and their
sample-level predictions, references, hit vectors, and token counts exactly.
The degenerate zero-width descriptive intervals reflect identical observed
outputs, not population-level certainty; with only three paired seeds and 20
samples per cell, we do not claim equivalence or non-inferiority.

\begin{figure}[t]
\centering
\includegraphics[width=\linewidth]{figures/vector_redesign/fig3_ppl_ruler.pdf}
\caption{\textbf{Perplexity stacking and RULER.} (a) PPL marginal change of
bf16 state under int4 KV (C4/PG19, 3 seeds). (b) RULER deltas (bf16 minus
fp32 state) for the five screen-selected cells (FWE and NIAH-multiquery,
4K/8K); thinking disabled, 20 samples per cell, and 3 paired dataset seeds.
Every observed delta is exactly zero in both attempts. Zero-width descriptive
intervals are not an equivalence test.}
\label{fig:ruler}
\end{figure}

\begin{figure}[t]
\centering
\includegraphics[width=\linewidth]{figures/vector_redesign/fig7_harness.pdf}
\caption{\textbf{Harness boundary and stacking cost.}
(a) Chunk-level write-back rounding: chunk=1 raises C4 PPL by roughly 87\%
relative to chunk=128 (1 seed, 1 sequence), so PPL is a chunk-level
approximation. (b) Marginal state-bf16 PPL delta under fp16 KV vs.\ int4 KV
(3 seeds): stacking introduces no measurable additional cost.}
\label{fig:harness}
\end{figure}

\begin{figure}[t]
\centering
\includegraphics[width=\linewidth]{figures/vector_redesign/fig8_gsm8k_per_seed.pdf}
\caption{\textbf{GSM8K per-seed paired accuracy.}
200 seeded items, greedy, no chain-of-thought, 9 dataset seeds. Thin lines
connect the same seed across allocations; the blue line and markers are seed
means. (a) 2B across four allocations; (b) 9B across two. This figure is the
per-seed resolution of Figure~\ref{fig:gsm8k}.}
\label{fig:gsm8k-seeds}
\end{figure}

\subsection{Per-layer sensitivity}
\label{sec:sensitivity}

We separately asked whether any single GDN layer should be kept at higher
precision. Bf16 was applied to one GDN layer at a time (18 configurations,
2 corpora, 3 seeds, 36 tests). Two raw intervals excluded zero (C4, layers 2
and 8, $p=0.049$ and $p=0.0036$), but neither survives Bonferroni
($\alpha/36$) or BH-FDR correction. There is no evidence for a per-layer
precision benefit, so the paper's allocation dimension is the whole-state
switch rather than layer-wise state precision
(Figure~\ref{fig:sensitivity}).

\begin{figure*}[t]
\centering
\includegraphics[width=\textwidth]{figures/vector_redesign/fig6_sensitivity.pdf}
\caption{\textbf{Per-layer state sensitivity (2B, C4/PG19, 3 seeds).}
Each row is bf16 applied to one GDN layer; points and paired 95\% CIs are
deltas against the fp32 baseline. Orange rings mark the two C4 layers with
raw $p<0.05$; neither is significant after Bonferroni or BH-FDR correction.}
\label{fig:sensitivity}
\end{figure*}

\subsection{Executable budget selection}
\label{sec:controller}

The controller turns the paper's two-dimensional accounting into an
executable decision. It selects \texttt{full}, \texttt{state\_only}, or
\texttt{joint} for the strict, medium, and high frozen budgets, respectively;
an infeasible request fails closed. Every selected mapping launches the
corresponding vLLM precision arguments and meets the requested 500 ms TTFT and
200 ms TPOT SLO in all three confirmatory seeds. In the independent temporal
run, mean goodput is 29.554, 29.552, and 38.905 req/s for strict, medium, and
high; mean P95 TTFT is 194.23, 192.28, and 372.42 ms, and mean P95 TPOT is
18.04, 17.84, and 38.77 ms. All nine parent-to-rerun metrics fall within the
predeclared environment-sensitive tolerances (10\% for goodput and 20\% for
tail latency), with 18,000/18,000 requests and no failures in either attempt.
This verifies policy execution for one Qwen3.5-2B, 4K, Random-workload,
single-GPU slice. It is not evidence that the selected allocation dominates
baselines across models or workloads.

\subsection{Serving}
\label{sec:serving}

Table~\ref{tab:serving} reports sustainable arrival rates for
\texttt{int4} and \texttt{int4+bf16} allocations under the v3 protocol, and
Figure~\ref{fig:serving} shows paired goodput deltas in the overload region.

\begin{table}[t]
\centering
\footnotesize
\caption{\textbf{Sustainable rate boundaries (req/s, 3 seeds, protocol v3).}
Formal matrix; values in parentheses are the second formal run (identical
contracts and seeds).
Differences between runs at Random60/1000 ms and ShareGPT/250 ms show that
single boundary cells are run-sensitive and are not claimed as effects.}
\label{tab:serving}
\setlength{\tabcolsep}{3.5pt}
\begin{tabular*}{\columnwidth}{@{\extracolsep{\fill}}l@{}lcccccc}
\toprule
Workload & Config & 250 & 500 & 1000 & 2000 & 3000\\
\midrule
Random60 & int4 & 30 & 35 & 35 (40) & 40 & 40\\
Random60 & int4+bf16 & 30 & 35 & 40 & 40 & 40\\
ShareGPT & int4 & 40 (35) & 40 & 40 & 40 & 40\\
ShareGPT & int4+bf16 & 35 (40) & 40 & 40 & 40 & 40\\
\bottomrule
\end{tabular*}
\end{table}

The only repeated boundary difference is Random60 at 1000 ms (35 vs.\ 40 in
the formal matrix), and it did not reappear in the second run (int4 also
reached 40 there). ShareGPT at 250 ms shows opposite directions between the
two runs. We therefore do not present either workload as improved; the
boundary tables are measurements, and we report them without assigning an
effect.

What is stable across the two runs is the paired goodput delta in the
Random60 overload region; because the second run uses identical contracts and
seeds, this establishes run-to-run stability rather than an independent
sample. At r40, bf16 state improves goodput by 0.334 [0.078, 0.589] at
250 ms (paired t-test $p=0.030$) and 0.215 [0.154, 0.276] at 500 ms
($p=0.004$); the second formal run measures 0.304 [0.250, 0.359]
($p=0.002$) and 0.138 [0.082, 0.195] ($p=0.009$). At r45/r50 and
2000--3000 ms the deltas are smaller but consistently positive in both runs
(e.g., r45/3000 ms: 0.367 [0.308, 0.426] vs.\ 0.372 [0.310, 0.435];
$p=0.001$ and $p=0.002$). None of the 60 serving cells survives BH-FDR at
$q<0.05$ (closest: formal r40/500 ms $q=0.066$; second run r40/250 ms
$q=0.052$). We therefore treat the overload-region deltas as directional
evidence whose weight comes from same-cell stability across the two
runs, and we stop short of treating them as sustainable-SLO improvements.
A cell-level direction audit over the same 60-cell grid is workload-limited:
all 13 Random60 cells that are overloaded in both runs have positive deltas in
both runs (13/13), while 7 of the 10 ShareGPT overload cells have opposite
signs, so the directional evidence does not generalize across workloads.
Per-cell $p$ and BH $q$ values for all 60 cells are archived with the
analysis artifacts. We report the goodput deltas as aggregate effects and
do not assign them to a mechanism. The controlled contrasts below delimit
which serving metrics reproduce.
ShareGPT at 500 ms and above shows boundary parity in both runs.

The mechanism matrix completes all 18 allocation-by-seed cells and 5,400
requests in each attempt. Throughput reproduces within 10\% in 18/18 cells
(maximum relative difference 5.10\%), but TTFT P95 does so in only 10/18
cells (maximum 75.39\%) and TPOT P95 in 17/18 (maximum 12.13\%). In the formal
attempt, joint-minus-full mean throughput is $-0.760$, $-0.815$, and
$-2.492$ req/s under fixed block count, fixed bytes, and fixed concurrency,
respectively; none of the nine metric-by-contrast tests survives Holm
correction. We therefore promote only the tested throughput direction and do
not attribute a universal causal mechanism or exact latency effect.

Finally, the four-configuration formal matrix completes 144/144 cells and
320,400/320,400 measurement requests with zero failures in each of two runs.
Nevertheless, its frozen Gate~4 fails: only 537/720 continuous goodput values
are within 10\% across attempts, 183/720 are outside tolerance, and the
maximum symmetric relative difference is 79.61\%. Sustainable labels are
more stable (713/720 cell-threshold labels and 38/40 all-seed boundaries
match), but these binary agreements do not override the failed primary gate.
We therefore use this matrix to disclose serving instability, not to claim a
four-configuration performance advantage.

\begin{figure*}[t]
\centering
\includegraphics[width=\textwidth]{figures/vector_redesign/fig4_serving.pdf}
\caption{\textbf{Serving, Random60 overload region.} (a) Paired goodput
delta (bf16 minus fp32 state) at rates 40--50 req/s, all thresholds, formal
and second run, 3 seeds each, 95\% CI. (b) bf16-minus-fp32 sustainable-rate
boundary changes for both workloads and runs. The three nonzero cells are
run-sensitive; none of the 60 cellwise comparisons survives BH-FDR.}
\label{fig:serving}
\end{figure*}

\section{Limitations}
\label{sec:limits}

\textbf{Harness boundary.} The perplexity harness rounds the state at chunk
boundaries (chunk 128). A chunk-1 ablation changes PPL from 19.35 to 36.16
($+87\%$), which shows the harness is not equivalent to per-token kernel
semantics. We treat PPL as supporting evidence and place the kernel-path
quality conclusions (GSM8K, RULER) on the vLLM serving path.

\textbf{RULER resolution.} The repaired protocol disables thinking and has
20 samples in each of five cells with three paired dataset seeds. Exact
agreement in both attempts is transparent but low-resolution evidence: it
does not define an equivalence margin or establish equivalence,
non-inferiority, or quality preservation outside the tested cells.

\textbf{Generalization.} All evidence is from Qwen3.5-2B/9B on one RTX 5090.
Claims are scoped to GDN-based hybrids; Mamba2-style architectures and other
hardware are untested. fp16-state quality is smoke-tested only; fp8/int8
state precisions are future work, so we make no claims about a full precision
spectrum.

\textbf{Serving variability.} Overload-region goodput is stable in direction
but not in magnitude across runs; boundary cells are run-sensitive. A
workload with very long requests or a different mix may not show the effect.
The broader four-configuration matrix does not pass its continuous-goodput
stability gate. We did not test offloading or prefix caching, each of which
changes the memory equation, and the four internal precision configurations
do not substitute for those external system baselines.

\textbf{Mechanism attribution.} We ran fixed-block-count, fixed-bytes, and
fixed-concurrency contrasts, but only throughput passes the temporal
reproducibility threshold in every cell. TTFT and one TPOT cell do not.
Moreover, the fixed-concurrency contrast changes the queueing constraint and
is not a pure dtype intervention. These results narrow the plausible
mechanism but do not identify a universal causal channel.

\textbf{Controller scope.} The selector is executable and reproducible, but
its verified evidence covers one Qwen3.5-2B, 4K, Random-workload, TP=1 slice
and three preregistered budget requests. It uses calibrated profiles and cold
restarts rather than online adaptation, and it does not interpolate to
unmeasured strata. Cross-model, cross-workload, cross-context, and switching
overhead claims require separate evaluation.

\textbf{Tensor parallelism.} The single-GPU ratios extend to tensor
parallelism only under uniform sharding. If the GDN state is sharded by its
hidden dimension and the attention KV by its KV heads, both $A\cdot L$ and
$G$ scale as $1/P$, so $r_{state}(L)$ is unchanged to first order; for
example, at $P=2$ with the 2B/int4 KV and $L=4096$ configuration, both terms
halve and the ratio is unchanged to first order. Per-rank
block rounding and communication overhead can shift the measured ratio, and
we did not test TP=2/4; the transfer is an expectation, not a measurement.

\textbf{Statistical resolution.} GSM8K at 2B detects a 1.0-point regression
with 67.5\% power; effects smaller than $\sim$1.16 points are below the
pre-registered MDE. Serving cells use $n=3$ seeds per cell and 60-cell BH-FDR
control; no cell survives $q<0.05$, so serving claims remain directional.
The GSM8K seed resolution is $n=9$ paired dataset seeds. The two serving runs
share contracts and seeds, so the second run provides run-stability evidence
rather than an independent sample. RULER uses only three paired dataset seeds;
its exact zero deltas are descriptive and do not supply an equivalence bound.
The M3 latency endpoints and M4 continuous-goodput family fail their frozen
reproducibility gates and remain unverified.

\section{Conclusion}
\label{sec:conclusion}

Hybrid linear-attention serving has two allocatable precision dimensions:
the attention-KV bit-width and the recurrent-state dtype. We measured the
capacity, quality, and serving consequences of treating both as a joint
budget and implemented a constraint-aware selector that produces executable
deployments. The verified result is the capacity effect: all 52 core pairs in
the 112-cell matrix favor bf16 state, with a 15.44\% median gain. The capacity
model is an idealized predictor with measured residuals, not a bound. Quality
remains task-specific: the 2B GSM8K regression is retained, while no-think
RULER shows exact observed agreement only in the five tested cells. The
controller's three budget mappings reproduce in one scoped slice. Controlled
serving contrasts reproduce throughput but not tail latency, and the broader
four-configuration matrix fails its continuous-goodput stability gate.

Operational rule. Under the target regime (short-context, high-concurrency,
memory-bound serving), the state switch is material: at 2B/4K/int4 KV it adds
roughly 247 concurrent sequence slots without changing the KV bytes
themselves. When the KV footprint dominates the budget (e.g., at longer
contexts), the relative gain shrinks, and operators should re-evaluate with
the same capacity and quality protocol before adopting the switch.

We release the full measurement contracts, per-seed data, and reproduction
artifacts so that the boundary cells and null results can be audited
independently. Long-form and retrieval evidence beyond GSM8K and RULER
(e.g., LongBench~\cite{longbench}) remains single-seed pilot evidence and is
excluded from the claims above.

\bibliographystyle{mlsys2025}
\bibliography{main}

\appendix
\section{Capacity model derivation}
\label{app:derivation}

For Qwen3.5-2B, the GDN state per sequence is derived from the vLLM
\texttt{MambaStateShapeCalculator} path: 18 layers $\times$ 1,085,440 bytes
per layer at fp32 (code-derived; the 9B state is 24.84 MiB). The per-token
attention-KV footprints $A_f$ and $A_q$ are computed from the GQA
head/layer configuration and the vLLM packed page layout (fp16: 2 bytes per
element; int4: 0.5 bytes per element). $G$ for bf16 is half of $G_{fp32}$.
Eq.~\eqref{eq:rstate} uses these derived values; the predicted ratios
reported in Table~\ref{tab:capacity} are evaluated with no fitted
parameters. The observed residual is attributable to the allocator's
discrete block rounding, which we verified from the per-attempt block-size
and block-count records.

For Eq.~\eqref{eq:capacity}, $N(L)$ is the number of concurrent length-$L$
sequences and $T(L)$ is their total token capacity.

\end{document}
